\documentclass[a4paper,11pt]{article}
\usepackage[french]{babel}
\usepackage{alphabeta}
\usepackage{amsmath}
\usepackage[osf,sups]{Baskervaldx} % lining figures
\usepackage[bigdelims,cmintegrals,vvarbb,frenchmath,baskervaldx]{newtxmath} % math font
\usepackage{graphicx}
\usepackage{xcolor}
\usepackage{url}
\usepackage[hidelinks]{hyperref}

% hyperrref configuration: remove ugly boxes and colors, keep urls in blue
\hypersetup{colorlinks,linkcolor=black,citecolor=black,urlcolor={blue!80!black}}


% paragraphs
\setlength{\parindent}{0em}
\setlength{\parskip}{1em}

% margins
\addtolength{\hoffset}{-4.5em}
\addtolength{\textwidth}{9em}
\addtolength{\voffset}{-4.5em}
\addtolength{\textheight}{9em}

% macros
\def\ud{\mathrm{d}}
\def\R{\mathbf{R}}

\begin{document}

{\bf Métriques sur un ouvert du plan}\\
Une métrique Riemanniene sur~$U\subseteq\R^2$ est la donnée de trois
fonctions~$E,F,G:U\mapsto\R$ telles que~$E>0$ et~$EG-F^2>0$.  On dit que la
métrique est conforme si~$E=G$ et~$F=0$.  Dans ce cas on écrit~$E=e^{2u}$.  La
métrique euclidienne est le cas particulier~$u=0$.  Les métriques conformes
s'écrivent ainsi en coordonées cartésienes et polaires
\[
	e^{2u}\left(\ud x^2+\ud y^2\right)
	=
	e^{2u}\left(\ud r^2+r^2\ud θ^2\right)
\]

{\bf Métrique induite par une paramétrisation}\\
Soit~$U\subseteq\R^2$ un ouvert du plan et~$φ:U\to\R^3$ une surface
paramétrée.  L'espace tangent à un point~$φ(x,y)$ de la surface est
généré par les vecteurs~$φ_x$ et~$φ_y$.  La matrice de la métrique
euclidienne sur cet espace tangent en cette base est la première forme
fondamentale
\[
	%\begin{pmatrix}E&F\\F&G\end{pmatrix}
	%=
	\begin{pmatrix}φ_x\cdotφ_x&φ_x\cdotφ_y\\
	φ_x\cdotφ_y&φ_y\cdotφ_y\end{pmatrix}
\]
Ces matrices déterminent une métrique Riemanniene sur~$U$ dite la métrique
induite par la paramétrisation.

{\bf Surface de révolution}\\
Une application~$φ:\R^2\mapsto\R^3$ de la forme
\[
	φ
	\begin{pmatrix}
		r\\θ
	\end{pmatrix}
	=
	\begin{pmatrix}
		R(r)\cosθ\\R(r)\sinθ\\z(r)
	\end{pmatrix}.
\]
La courbe génératrice est~$r\mapsto\begin{pmatrix}R(r)\\z(r)\end{pmatrix}$.
La métrique induite est donnée par~$E=(R')^2+(z')^2$, $F=0$, $G=R^2$.

{\bf Plongement isométrique local}\\
Étant donnée une métrique sur~$U\subseteq\R^2$, on peut se poser la question
si elle est la métrique induite par une paramétrisation~$φ:U\to\R^3$.
Cela équivaut à résoudre le système de trois EDP d'ordre~$1$ suivant:
\[
	\begin{cases}
		E=φ_x\cdotφ_x\\
		F=φ_x\cdotφ_y\\
		G=φ_y\cdotφ_y\\
	\end{cases}
\]

{\bf Plongement d'une métrique conforme une surface de révolution}\\
Si on veut plonger une métrique avec facteur conforme radial~$r\mapsto
e^{2u(r)}$, c'est naturel de penser à une surface de révolution.
L'équation à résoudre est
un cas particulier du système d'EDP antérieur
\[
	\begin{cases}
		(R')^2+(z')^2=e^{2u}\\
		R^2=r^2e^{2u}
	\end{cases}
\]
qui admet une solution explicite.  On trouve par la deuxième équation
\[
	R(r)=re^{u(r)}
\]
de sorte que~$R'(r)=(1+ru'(r))e^{u(r)}$.
En remplaçant dans l'autre équation~:
\[
	(z'(r))^2=e^{2u(r)}\sqrt{1-\left(1+ru'(r)\right)^2}
\]
cette équation aura solution quand~$(1+ru')^2\le 1$, ou équivalemment
\begin{equation}\label{eq:compatirev}
	ru'(r)\in[-2,0]
\end{equation}
On appelle la condition~\ref{eq:compatirev} la ``condition de révolution''.
Les fonctions~$u$ qui satisfont cette condition déterminent des métriques
conformes qui peuvent se plonger isométriquement dans une surface de
révolution.

{\bf Jacobi-Maupertuis}\\
Les trajectoires d'un système conservatif de potentiel~$V$ sont les solutions
de l'équation~$\ddot q=-\nabla V(q)$.  L'énérgie d'une trajectoire
est la quantité conservée~$E=\tfrac12\|\dot q\|^2+V(q)$.  Principe de
Maupertuis~: les trajectoires d'énergie totale~$E$ sont, à réparamétrisation
près, des géodésiques de la métrique de Jacobi~$\ud s^2=2(E-V(q)\ud q^2$.

{\bf Potentiels centraux homogènes}\\
On considère des potentiels de la forme~$V(r)=r^α/α$.  Pour~$E=0$ le facteur
conforme de Jacobi est~$e^{2u(r)}=-2r^α/α$, ce qui force~$α<0$.  On a
donc~$u(r)=\frac{α}{2}\log(r)+c$ et~$u'(r)=\frac{α}{2r}$. La condition de
révolution pour ces métriques est alors~$α\in[-4,0]$.







\end{document}


% vim:set tw=79 sw=2 ts=2:
